Tomando siempre el cociente «ensayo 2 entre ensayo 1», la energía disipada es cuatro veces mayor durante el mismo intervalo \([0,2\pi/\omega]\), pero sólo dos veces mayor al comparar un periodo propio de cada respuesta:
\[
D_1=\pi cA^2\omega,\qquad
D_{2,\,\text{mismo tiempo}}=4\pi cA^2\omega,
\qquad D_{2,\,\text{un periodo}}=2\pi cA^2\omega.
\]
El trabajo del amortiguador es \(-D\); la excitación externa aporta un trabajo neto \(+D\) en cada intervalo. La respuesta rápida disipa más por unidad de tiempo, pero su periodo dura la mitad.
